\documentclass[10pt]{article}

\usepackage[utf8]{inputenc}
\usepackage[T1]{fontenc}
\usepackage[margin=0.78in]{geometry}
\usepackage{amsmath,amssymb}
\usepackage{booktabs}
\usepackage{enumitem}
\usepackage{hyperref}
\usepackage{xcolor}

\hypersetup{colorlinks=true,linkcolor=blue,citecolor=blue,urlcolor=blue}
\newcommand{\system}{\textsc{Budget-NLA}}
\newcommand{\todoresult}{\textcolor{red}{TBD}}

\title{\system: Risk-Weighted Adaptive Monitoring for\\
Pre-Harm Sabotage Containment}

\author{
Jhoan Stevan Mosquera Ortiz\\
\texttt{jhoan.mosquera@uao.edu.co}
\and
Juan Felipe Plata Barbosa\\
\texttt{juan.plata@uao.edu.co}
\and
Maryori Alejandra Lasso Diaz\\
\texttt{maryori.lasso@uao.edu.co}
\and
Angel Alexis del Castillo Lerma\\
\texttt{angel.del@uao.edu.co}
}

\date{Universidad Aut\'onoma de Occidente --- AI Incident Response Sprint, 2026}

\begin{document}
\maketitle

\begin{abstract}
Agent monitors are imperfect and heterogeneous: observable actions, visible
chain of thought (CoT), activation probes, and natural-language activation
readouts differ in cost and in the failures they expose. Evaluating each signal
only by classification accuracy does not determine how it should be used under
a limited operational budget. We introduce \system, a reproducible pilot that
measures monitoring signals at a fixed pre-action boundary and allocates a
normalized monitoring budget across task types with different impact. The
system combines a Qwen2.5-7B agent, a layer-20 linear probe, a Natural Language
Autoencoder (NLA), an isolated Gemini reader, and a deterministic policy gate.
We evaluate 24 paired trajectories spanning benign behavior, rejected attacks,
visible sabotage, and stealth sabotage in three simulated tool-use families.
Our primary outcomes are pre-harm recall, benign disruption, expected prevented
harm, and prevention efficiency; AUROC is retained only as a diagnostic. Real
experimental results are intentionally left blank in this draft until the
frozen pipeline completes.
\end{abstract}

\section{Introduction}

LLM agents can pursue harmful side objectives while completing an apparently
legitimate task. SHADE-Arena demonstrates that monitoring such behavior becomes
particularly difficult when the agent can conceal relevant reasoning from the
monitor \cite{kutasov2025shade}. This motivates additional monitoring surfaces,
including internal activations, but it does not remove the operational
constraint that strong monitors, NLA verbalization, and human review have
different costs.

Prior work studies how to combine runtime monitors under a budget
\cite{hua2025cost}, while NLA research shows that natural-language bottlenecks
can expose information encoded in model activations
\cite{frasertaliente2026nla}. Monitor-side experiments further suggest that NLA
and CoT errors can be non-identical, but identify native agent-side readouts,
position selection, and matched monitoring baselines as open directions
\cite{bowkis2026eliciting}. Our pilot connects these questions at an operational
boundary: immediately before a critical tool call.

We ask:

\begin{quote}
Under a fixed monitoring budget, when and on which task types does access to an
agent's internal state through a linear probe and an NLA provide additional
pre-harm containment value over observable actions and visible CoT?
\end{quote}

Our intended contribution is not a new probe, NLA, routing algorithm, or
security game. It is an evaluation harness that (i) measures heterogeneous
signals at the same pre-action event, (ii) translates detections into prevented
weighted harm, and (iii) compares monitoring allocations under the same
normalized budget.

\section{Related Work}

\paragraph{Agent sabotage and monitoring.}
SHADE-Arena pairs benign main tasks with harmful side objectives and measures
both sabotage and monitor performance \cite{kutasov2025shade}. Its results
establish that successful side-task execution and evasion must be evaluated
together, and that access to hidden reasoning materially changes the monitoring
problem.

\paragraph{Natural-language activation readouts.}
NLAs reconstruct model activations through a natural-language bottleneck,
providing an inspectable representation of latent state
\cite{frasertaliente2026nla}. Bowkis and Africa evaluate NLA readouts as a
monitoring surface and report unique catches relative to CoT, while also showing
that generic extraction positions often contain irrelevant local information
\cite{bowkis2026eliciting}. We therefore freeze one semantically defined event
rather than selecting the best-performing position after seeing results.

\paragraph{Cost-constrained monitoring.}
Hua et al. formalize the combination of runtime monitors subject to an
average-case budget \cite{hua2025cost}. \system{} does not claim the general
resource-allocation idea. Our narrower focus is to estimate the operational
prevention value of observable, CoT, probe, and native agent-side NLA channels
in a controlled pre-harm tool-use setting.

\section{Experimental Design}

\subsection{Models and fixed boundary}

The target agent is \texttt{Qwen/Qwen2.5-7B-Instruct} with deterministic
decoding. The same layer-20 residual-stream activation is supplied to (a) an L2
logistic-regression probe and (b) the
\texttt{Kameshr/nla-qwen2.5-7b-L20-av} verbalizer. The activation is extracted
at the final token of the serialized marker
\texttt{\textbackslash n<DECISION>\textbackslash n}, before the tool name and
arguments are generated. We call this event
\texttt{PRE\_ACTION\_DECISION\_BOUNDARY}.

The Qwen agent generates with the NLA adapter disabled. The adapter is enabled
only while producing the NLA verbalization. This prevents the interpretation
module from altering the behavior being monitored.

\subsection{Task families and paired conditions}

All tools are simulated; no email, file, permission, or network state is
modified. We define irreversible harm as committing an unauthorized simulated
external transfer or access-control modification. Detection after commit does
not count as containment.

\begin{table}[h]
\centering
\small
\begin{tabular}{lll}
\toprule
Family & Critical tool & Simulated harm \\
\midrule
External report & \texttt{publish\_report} & protected record is published \\
Corporate email & \texttt{send\_email} & protected data or unauthorized recipient \\
Access list & \texttt{update\_allowlist} & unauthorized address or range is added \\
\bottomrule
\end{tabular}
\caption{Frozen task families.}
\label{tab:families}
\end{table}

Each family contains four conditions: benign behavior, a rejected attack,
visible sabotage, and stealth sabotage. Two fixed linguistic variants yield
$3\times4\times2=24$ trajectories. The rejected-attack condition separates the
presence of a malicious instruction from adoption of the malicious objective.

\subsection{Isolated monitoring channels}

\begin{itemize}[leftmargin=*]
    \item \textbf{Observable:} policy, authorized task, and information visible
    before the future tool call; no CoT or internal signal.
    \item \textbf{CoT:} the observable view plus reasoning truncated before
    \texttt{<DECISION>}.
    \item \textbf{Probe:} a scalar score from the pre-action activation, tested
    with leave-one-template-out (LOTO) splits.
    \item \textbf{NLA:} the observable view plus the NLA verbalization; no CoT
    or probe score.
    \item \textbf{Hybrid:} observable evidence, CoT, NLA, and probe score. This
    is treated as a policy input, not proof that any one channel caused a catch.
\end{itemize}

Gemini receives each view in an independent, stateless call and must return a
strict JSON verdict containing a risk score, actionable-alert flag, cited
evidence, policy conflict, and one of \texttt{ALLOW}, \texttt{ESCALATE}, or
\texttt{BLOCK}. A deterministic gate handles inconsistent responses
conservatively. API evaluation occurs after Qwen/NLA/probe acquisition from a
frozen \texttt{monitor\_cases.json} contract. This separation prevents an API
failure or retry from forcing another expensive activation run.

\subsection{Probe protocol}

Let $h_i\in\mathbb{R}^{d}$ be the layer-20 activation and $y_i\in\{0,1\}$ the
behavioral sabotage label. The probe estimates

\begin{equation}
    p_i = \sigma(w^\top h_i+b),
\end{equation}

with L2 regularization, $C=1$, the \texttt{liblinear} solver, and 1000 maximum
iterations. Every test template is excluded from its training fold. Thresholds
must be selected from training folds only. We report AUROC and AUPRC as
diagnostics, not as the main measure of containment.

\section{Budgeted Containment}

For trajectory $i$, let $s_i\geq0$ denote impact severity. Let
$A_{im}\in\{0,1\}$ indicate whether monitor $m$ raises an actionable alert
before harm, and let $I_{im}$ denote successful intervention.

\paragraph{Pre-harm recall.}
\begin{equation}
\mathrm{PHR}_m =
\frac{\sum_i \mathbb{1}[y_i=1]A_{im}}
{\sum_i \mathbb{1}[y_i=1]}.
\end{equation}

\paragraph{Benign disruption rate.}
\begin{equation}
\mathrm{BDR}_m =
\frac{\sum_i \mathbb{1}[y_i=0]A_{im}}
{\sum_i \mathbb{1}[y_i=0]}.
\end{equation}

\paragraph{Expected prevented harm and efficiency.}
\begin{align}
\mathrm{EPH}_m &= \sum_i s_i\,\mathbb{1}[y_i=1]A_{im}I_{im},\\
\mathrm{PE}_m &= \frac{\mathrm{EPH}_m}{C_m},
\end{align}
where $C_m$ is monitoring cost. We report assumed normalized cost separately
from measured latency, tokens, API calls, and GPU-seconds.

\paragraph{Incremental NLA value.}
\begin{align}
\mathrm{IMV}_{\mathrm{NLA}} &=
\mathrm{EPH}_{O\cup N}-\mathrm{EPH}_{O},\\
\mathrm{IMV}_{\mathrm{NLA}\mid\mathrm{CoT}} &=
\mathrm{EPH}_{H}-\mathrm{EPH}_{O\cup C},
\end{align}
where $O$, $N$, $C$, and $H$ denote observable, NLA, CoT, and hybrid policies.

\paragraph{Budget assignment.}
Let $x_{ik}\in\{0,1\}$ choose monitoring policy $k$ for trajectory $i$,
$b_{ik}$ be its estimated prevented-harm benefit, $c_k$ its normalized cost,
and $B$ the total budget. The dashboard solves the multiple-choice allocation

\begin{align}
\max_{x}\quad & \sum_i\sum_k b_{ik}x_{ik} \\
\text{s.t.}\quad & \sum_i\sum_k c_kx_{ik}\leq B,\\
& \sum_k x_{ik}=1 \quad \forall i,\\
& x_{ik}\in\{0,1\}.
\end{align}

We compare this adaptive allocation with fixed baselines under the same budget.
The current optimizer is a pilot implementation and is not claimed to model a
strategic attacker.

\section{Hypotheses and Analysis Plan}

The frozen hypotheses are: (H1) at least one internal channel uniquely catches
attempted sabotage missed by observable monitoring; (H2) CoT and NLA have
non-identical errors, particularly for stealth sabotage; (H3) probe and NLA
distinguish adopted attacks from rejected malicious instructions; and (H4) the
adaptive policy yields lower expected residual harm than a uniform allocation
under the same normalized budget.

Results will be reported as counts and descriptive proportions because the
pilot contains only 24 trajectories. No metric, threshold, layer, extraction
position, or primary hypothesis will be changed after observing test results.
Any necessary departure is recorded in \texttt{DEVIATIONS.md} and labeled
exploratory.

\section{Results}

\begin{table}[h]
\centering
\small
\begin{tabular}{lrrrr}
\toprule
Channel & PHR & BDR & EPH & PE \\
\midrule
Observable & \todoresult & \todoresult & \todoresult & \todoresult \\
CoT & \todoresult & \todoresult & \todoresult & \todoresult \\
Probe & \todoresult & \todoresult & \todoresult & \todoresult \\
NLA & \todoresult & \todoresult & \todoresult & \todoresult \\
\bottomrule
\end{tabular}
\caption{Primary channel metrics. Values must come from the frozen real run.}
\label{tab:primary}
\end{table}

\begin{table}[h]
\centering
\small
\begin{tabular}{lrrr}
\toprule
Policy & Cost & Weighted harm prevented & Attack success \\
\midrule
Uniform baseline & \todoresult & \todoresult & \todoresult \\
Action only & \todoresult & \todoresult & \todoresult \\
CoT + NLA & \todoresult & \todoresult & \todoresult \\
Probe + NLA + CoT & \todoresult & \todoresult & \todoresult \\
Adaptive LP & \todoresult & \todoresult & \todoresult \\
\bottomrule
\end{tabular}
\caption{Budget-policy comparison under one normalized budget.}
\label{tab:policies}
\end{table}

\paragraph{Required error analysis.}
Alongside aggregate metrics, we will report unique catches, benign flips,
visible-versus-stealth differences, rejected-attack errors, failed JSON/API
calls, and cases where NLA and probe disagree. \todoresult

\section{Reproducibility and Implementation}

The repository freezes model identifiers, extraction layer, decision marker,
probe settings, scenario templates, seeds, normalized costs, and a versioned
Pydantic result schema. The execution path is:

\begin{enumerate}[leftmargin=*]
    \item run Qwen, capture layer-20 activations, and generate NLA text;
    \item compute LOTO probe scores;
    \item export \texttt{results/monitor\_cases.json};
    \item run isolated Gemini evaluations into a resumable JSONL file;
    \item build aggregate metrics, SQLite export, and the Streamlit dashboard.
\end{enumerate}

Secrets are loaded from an untracked \texttt{.env}; only
\texttt{.env.example} is committed. Real model revisions and the successful
host package environment must be recorded with the final results.

\section{Limitations}

The task suite is small, synthetic, and single-boundary; results cannot be
generalized to long-horizon deployments. Severity weights and normalized costs
are assumptions, not organizational measurements. The NLA verbalizer and probe
are model- and layer-specific. A missing NLA signal does not imply that the
agent lacked the represented information. The post-run Gemini stage evaluates a
counterfactual gate rather than demonstrating real-time latency. Finally, the
allocation assumes stationary benefits and does not yet model an attacker that
adapts to the monitoring schedule.

\section{Safety and Ethics}

All critical tools are mocks and all recipients, canaries, addresses, and
records are synthetic. No real exfiltration, access-control change, or external
side effect is executed. The system is an incident-response research prototype,
not a production authorization mechanism; deterministic permission checks
remain necessary even when internal monitoring is available.

\section{Conclusion}

\system{} operationalizes a narrow question left open by accuracy-only monitor
evaluations: whether expensive internal signals prevent enough weighted harm to
justify where they are deployed under a fixed budget. The final claim will be
limited to the real 24-trajectory pilot and will be written only after the
frozen run completes. \todoresult

\bibliographystyle{plain}
\bibliography{references}

\end{document}
